% ═══════════════════════════════════════════════════════════════
% AB-02a: La familia -- Grundwortschatz
% Abgeleitet aus LE-02a (Score 9.3/10)
% ═══════════════════════════════════════════════════════════════

\documentclass[11pt, a4paper]{article}

% ═══════════════════════════════════════════════════════════════
% SW-MODUS TOGGLE
% ═══════════════════════════════════════════════════════════════
% Setze \swmodustrue für Schwarz-Weiß-Druck
\newif\ifswmodus
\swmodusfalse    % Standard: Farbe

\usepackage[utf8]{inputenc}
\usepackage[T1]{fontenc}
\usepackage[ngerman]{babel}
\usepackage{helvet}
\renewcommand{\familydefault}{\sfdefault}

\usepackage[margin=2cm]{geometry}
\usepackage{enumitem}
\usepackage{tabularx}
\usepackage{booktabs}
\usepackage{tikz}

\usepackage{xcolor}
\usepackage{tcolorbox}
\tcbuselibrary{skins}

\usepackage{fancyhdr}
\usepackage{lastpage}
\usepackage{needspace}

% Farben - Basis
\definecolor{spanisch-color}{HTML}{c41e3a}
\definecolor{grau-color}{HTML}{888888}
\definecolor{hellgrau-color}{HTML}{f5f5f5}
\definecolor{orange-color}{HTML}{b35c00}

% SW-Farben (druckerfreundlich)
\definecolor{sw-dark}{HTML}{000000}
\definecolor{sw-gray}{HTML}{666666}
\definecolor{sw-white}{HTML}{ffffff}

% Bedingte Farbzuweisung
\ifswmodus
    \colorlet{spanisch}{sw-dark}
    \colorlet{grau}{sw-gray}
    \colorlet{hellgrau}{sw-white}
    \colorlet{orange}{sw-dark}
\else
    \colorlet{spanisch}{spanisch-color}
    \colorlet{grau}{grau-color}
    \colorlet{hellgrau}{hellgrau-color}
    \colorlet{orange}{orange-color}
\fi

\newcommand{\fachfarbe}{spanisch}

% Variablen
\newcommand{\thema}{La familia -- Grundwortschatz}
\newcommand{\sequenz}{02}
\newcommand{\block}{a}
\newcommand{\fach}{Spanisch}
\newcommand{\klasse}{7}
\newcommand{\maxpunkte}{20}

% Kopf-/Fußzeile
\pagestyle{fancy}
\fancyhf{}
\renewcommand{\headrulewidth}{0.4pt}
\renewcommand{\footrulewidth}{0.4pt}
\lhead{\fach{} -- Klasse \klasse}
\chead{AB-\sequenz\block}
\rhead{}
\lfoot{}
\rfoot{Seite \thepage\ von \pageref{LastPage}}

% Befehle
\newcommand{\punktebox}[1]{\hfill\fbox{\small \_/#1}}
\newcommand{\luecke}[1]{\rule{#1}{0.4pt}}
\newcommand{\es}[1]{\textcolor{\fachfarbe}{\textbf{#1}}}

% Merkkasten (SW-kompatibel: schwebender Titel)
\newtcolorbox{merkkasten}[1][]{
    colback=hellgrau,
    colframe=\fachfarbe,
    colbacktitle=white,
    coltitle=\fachfarbe,
    boxrule=1pt,
    arc=0pt,
    left=8pt, right=8pt, top=8pt, bottom=8pt,
    fonttitle=\bfseries,
    attach boxed title to top left={yshift=-2mm, xshift=5mm},
    boxed title style={boxrule=0pt, colback=white},
    title=#1
}

% Hinweis (SW-kompatibel: schwebender Titel)
\newtcolorbox{hinweis}{
    colback=white,
    colframe=orange,
    colbacktitle=white,
    coltitle=orange,
    boxrule=1pt,
    arc=0pt,
    left=5pt, right=5pt, top=5pt, bottom=5pt,
    fonttitle=\bfseries,
    attach boxed title to top left={yshift=-2mm, xshift=5mm},
    boxed title style={boxrule=0pt, colback=white},
    title=Hinweis
}

% Vokabelkasten
\newtcolorbox{vokabelkasten}{
    colback=white,
    colframe=\fachfarbe,
    boxrule=0.5pt,
    arc=0pt,
    left=5pt, right=5pt, top=5pt, bottom=5pt
}

% Aufgabe (kein Seitenumbruch innerhalb)
\newenvironment{aufgabe}[1]{%
    \needspace{5\baselineskip}%
    \nopagebreak[4]%
    \vspace{10pt}
    \noindent\textbf{\textcolor{\fachfarbe}{Aufgabe #1}}
    \nopagebreak[4]%
    \vspace{5pt}
    \nopagebreak[4]%
    \begin{list}{}{\leftmargin=0pt}
    \item[]
}{%
    \end{list}
}

% ═══════════════════════════════════════════════════════════════
\begin{document}

% Kopfbereich
\begin{center}
    {\Large\bfseries\textcolor{\fachfarbe}{Arbeitsblatt: \thema}}
\end{center}

\vspace{5pt}

\noindent
\begin{tabularx}{\textwidth}{|X|X|X|}
\hline
\textbf{Name:} \luecke{4cm} &
\textbf{Klasse:} \klasse &
\textbf{Datum:} \luecke{2cm} \\
\hline
\end{tabularx}

\vspace{15pt}

% ═══════════════════════════════════════════════════════════════
% MERKKASTEN: Familienmitglieder
% ═══════════════════════════════════════════════════════════════

\begin{merkkasten}[Merkkasten: Die Familie -- La familia]

\begin{center}
\begin{tabular}{|l|l||l|l|}
\hline
\textbf{Spanisch} & \textbf{Deutsch} & \textbf{Spanisch} & \textbf{Deutsch} \\
\hline
\es{el padre} & \luecke{2.5cm} & \es{la madre} & \luecke{2.5cm} \\
\hline
\es{el hermano} & \luecke{2.5cm} & \es{la hermana} & \luecke{2.5cm} \\
\hline
\es{el abuelo} & \luecke{2.5cm} & \es{la abuela} & \luecke{2.5cm} \\
\hline
\es{el tío} & \luecke{2.5cm} & \es{la tía} & \luecke{2.5cm} \\
\hline
\es{el primo} & \luecke{2.5cm} & \es{la prima} & \luecke{2.5cm} \\
\hline
\end{tabular}
\end{center}

\vspace{0.5em}
\textit{Übertrage die Übersetzungen aus dem Lernpfad oder von der Tafel!}
\end{merkkasten}

\vspace{10pt}

% ═══════════════════════════════════════════════════════════════
% AUFGABE 1: Zuordnung (4P)
% ═══════════════════════════════════════════════════════════════

\begin{aufgabe}{1}
Verbinde die spanischen Wörter mit der deutschen Übersetzung. \punktebox{4}
\end{aufgabe}

\begin{center}
\begin{tabular}{rcl}
\es{el padre} & \luecke{3cm} & die Schwester \\[0.8em]
\es{la madre} & \luecke{3cm} & der Bruder \\[0.8em]
\es{el hermano} & \luecke{3cm} & die Mutter \\[0.8em]
\es{la hermana} & \luecke{3cm} & der Vater \\
\end{tabular}
\end{center}

\vspace{15pt}

% ═══════════════════════════════════════════════════════════════
% AUFGABE 2: Lückentext (4P)
% ═══════════════════════════════════════════════════════════════

\begin{aufgabe}{2}
Ergänze die Lücken mit dem passenden Familienmitglied. \punktebox{4}
\end{aufgabe}

\begin{vokabelkasten}
\textit{Wortliste:} padre | madre | hermano | hermana | abuelo
\end{vokabelkasten}

\vspace{0.5em}

\noindent
a) Mi \luecke{3cm} se llama Carlos. Er ist mein Vater. \\[0.8em]
b) Mi \luecke{3cm} se llama María. Sie ist meine Mutter. \\[0.8em]
c) Tengo un \luecke{3cm} y una \luecke{3cm}. (Bruder und Schwester) \\[0.8em]
d) Mi \luecke{3cm} es muy simpático. Er ist der Vater meines Vaters. \\

\vspace{15pt}

% ═══════════════════════════════════════════════════════════════
% AUFGABE 3: Stammbaum beschriften (6P)
% ═══════════════════════════════════════════════════════════════

\begin{aufgabe}{3}
Beschrifte den Stammbaum mit den richtigen Familienmitgliedern. \punktebox{6}
\end{aufgabe}

\begin{hinweis}
Verwende: el abuelo, la abuela, el padre, la madre, el tío, la tía, el hermano / la hermana, el primo, la prima
\end{hinweis}

\vspace{0.5em}

\begin{center}
\begin{tikzpicture}[
    every node/.style={draw, minimum width=2.2cm, minimum height=0.7cm, align=center},
    level 1/.style={sibling distance=6cm},
    level 2/.style={sibling distance=3cm},
    edge from parent/.style={draw, -}
]

% Großeltern
\node (abuelo) at (0,4) {\small 1. \luecke{1.5cm}};
\node (abuela) at (3,4) {\small 2. \luecke{1.5cm}};

% Linie zwischen Großeltern
\draw (abuelo.east) -- (abuela.west);

% Elterngeneration
\node (padre) at (-1.5,2) {\small 3. \luecke{1.5cm}};
\node (madre) at (1.5,2) {\small 4. \luecke{1.5cm}};
\node (tio) at (4.5,2) {\small 5. \luecke{1.5cm}};
\node (tia) at (7.5,2) {\small 6. \luecke{1.5cm}};

% Linien von Großeltern zu Eltern
\draw (1.5,3.6) -- (1.5,3) -- (-1.5,3) -- (padre.north);
\draw (1.5,3) -- (1.5,2.7) -- (madre.north);
\draw (1.5,3.6) -- (1.5,3) -- (4.5,3) -- (tio.north);
\draw (4.5,3) -- (7.5,3) -- (tia.north);

% Linien zwischen Ehepaaren
\draw (padre.east) -- (madre.west);
\draw (tio.east) -- (tia.west);

% Kindergeneration
\node (yo) at (-1.5,0) {\small 7. \textbf{YO}};
\node (hermano) at (1.5,0) {\small 8. \luecke{1.5cm}};
\node (primo) at (4.5,0) {\small 9. \luecke{1.5cm}};
\node (prima) at (7.5,0) {\small 10. \luecke{1.5cm}};

% Linien von Eltern zu Kindern
\draw (0,1.3) -- (0,0.7) -- (-1.5,0.7) -- (yo.north);
\draw (0,0.7) -- (1.5,0.7) -- (hermano.north);
\draw (6,1.3) -- (6,0.7) -- (4.5,0.7) -- (primo.north);
\draw (6,0.7) -- (7.5,0.7) -- (prima.north);

\end{tikzpicture}
\end{center}

\vspace{15pt}

% ═══════════════════════════════════════════════════════════════
% AUFGABE 4: Mini-Dialog (6P)
% ═══════════════════════════════════════════════════════════════

\begin{aufgabe}{4}
Schreibe einen kurzen Dialog. Nutze die Redemittel. \punktebox{6}
\end{aufgabe}

\begin{hinweis}
\textbf{Redemittel:} ¿Quién es? | Es mi... | Se llama... | ¿Y quién es él/ella?
\end{hinweis}

\vspace{0.5em}

\noindent
\textbf{A:} (zeigt auf ein Foto) \es{¿Quién es?}

\vspace{0.3em}
\textbf{B:} \luecke{8cm}

\vspace{0.3em}
\textbf{A:} \es{¿Y quién es ella?} (zeigt auf eine Frau)

\vspace{0.3em}
\textbf{B:} \luecke{8cm}

\vspace{0.3em}
\textbf{A:} \luecke{8cm}

\vspace{0.3em}
\textbf{B:} \luecke{8cm}

\vspace{20pt}
\hrule
\vspace{10pt}

\begin{flushright}
\textbf{Punkte:} \luecke{1cm} / \maxpunkte
\end{flushright}

\end{document}
